\documentclass[12pt]{article}
\usepackage[letterpaper,top=2cm,bottom=2cm,left=3cm,right=3cm]{geometry}
\usepackage[utf8]{inputenc}
\usepackage[T1]{fontenc}
\usepackage{amsmath,amssymb,amsthm,mathtools}
\usepackage{bbm}
\usepackage{booktabs}
\usepackage[colorlinks,citecolor=blue,urlcolor=blue]{hyperref}
\usepackage{cleveref}

\newcommand*{\eg}{\emph{e.g.}{}}
\newcommand*{\ie}{\emph{i.e.}{}}
\newcommand{\norm}[1]{\left\lVert#1\right\rVert}

\theoremstyle{plain}
\newtheorem{theorem}{Theorem}[section]
\newtheorem{lemma}[theorem]{Lemma}
\newtheorem{assumption}{Assumption}
\newtheorem{prop}{Proposition}
\newtheorem{definition}[theorem]{Definition}
\newtheorem{remark}{Remark}

\title{Working Formalization:\\
Stable Post-Hoc Inference for Optimization-Selected Policies}
\author{Nicholas Mino \and Wenbin Zhou \and Shixiang Zhu\\
Carnegie Mellon University}
\date{}

\begin{document}
\maketitle

\begin{abstract}
This note records the working problem formulation for the project
\emph{Algorithmic Stability for Post-Hoc Selection}.
It expands the end formalization from the project draft:
stability of constrained, data-driven optimizations, and post-hoc
inference after stable selection.
It is a formulation document, not a theorem paper.
\end{abstract}

\section{Introduction}
We observe data, run a data-dependent optimization-based selection
algorithm, and then want to make an inference or evaluation statement
about the selected object.
The difficulty is that selection uses the same data as inference, so the
selected object is statistically dependent on the data used to assess it.

The goal is to reduce or avoid data splitting by finding a general
framework for post-hoc optimization-based selection algorithms:
convert an optimization-based selection algorithm into a
\emph{stable} selection algorithm.
Such a framework should identify when a selection rule can be randomized
or otherwise controlled so that its output is sufficiently and
certifiably stable for valid downstream inference or evaluation.

The post-hoc inference problem is itself very general.
We narrow it by focusing on a concrete objective---uncertainty
quantification---in which a hyperparameter is chosen by constrained
optimization on shared calibration data.
The overall objective mixes algorithm stability, post-hoc inference,
and optimization theory, and is decomposed into two parts:
\begin{enumerate}
\item stability of constrained optimizations;
\item post-hoc inference on top of those stability results.
\end{enumerate}

\section{Common Setup}
Let
\[
D=(Z_1,\dots,Z_n)
\]
denote the observed dataset, where
\[
Z_1,\dots,Z_n \overset{\mathrm{i.i.d.}}{\sim} P
\]
and \(P\) is an unknown population distribution.

Let \(A\) be a data-dependent optimization-based selection algorithm, and let
\[
\widehat S = A(D)
\]
denote the selected object.

We seek a stabilized selection rule \(\widetilde A\) producing
\[
\widetilde S = \widetilde A(D;\xi),
\]
where \(\xi\) represents additional randomness or another stabilizing
device.
Depending on the mechanism, \(\xi\) may represent noise, subsampling
randomness, randomized optimization perturbation, privacy randomness,
or another source of algorithmic randomization.

The target post-selection inference guarantee is
\[
\mathbb{P}\bigl(\theta_{\widetilde S}\in C_{\widetilde S}(D)\bigr)
\ge 1-\alpha.
\]
In words: after using the data to select an object, the confidence set
reported for that selected object should contain its true parameter with
probability at least \(1-\alpha\).

\paragraph{Policy view.}
We treat the selected object as (or as inducing) a \emph{policy}:
the solution of an optimization problem that maps context to a decision
or hyperparameter.
Stability is formalized for the policy induced by that optimization.
The stabilization mechanism is then chosen so that the induced policy
is stable enough for the target guarantee above.

\section{Narrowed Setting: Uncertainty Quantification via Constrained Optimization}
Consider a post-hoc inference setting in uncertainty quantification,
where a hyperparameter \(\theta\) is determined by optimization using
shared calibration data.

Let \(\mathcal{C}(X;\mathcal{D},\theta)\) denote an uncertainty set
constructed from dataset \(\mathcal{D}\) and hyperparameter \(\theta\).
Assume that, for any fixed \(\mathcal{D}\) and any prespecified \(\theta\),
the set satisfies \emph{validity}:
\begin{equation}
\label{eq:validity}
\mathbb{P}\bigl(Y\in\mathcal{C}(X;\mathcal{D},\theta)\bigr)
\ge 1-\alpha.
\end{equation}

Suppose there exist two data-dependent functions \(\hat f\) and \(\hat g\)
(with \(\sigma(\hat f)\subseteq\sigma(\mathcal{D})\) and
\(\sigma(\hat g)\subseteq\sigma(\mathcal{D})\)) that jointly define a
constrained optimization problem determining \(\theta\):
\begin{equation}
\label{eq:opt}
\pi:\qquad
\min_{\theta}\ \hat f(y;\theta)
\quad\text{s.t.}\quad
\hat g(\theta)\le 0.
\end{equation}
Let \(\hat\theta(\mathcal{D})\) denote the output of~\eqref{eq:opt}.

The questions are:
\begin{enumerate}
\item[(i)] under what conditions does~\eqref{eq:validity} still hold when
\(\hat\theta(\mathcal{D})\) is plugged in for \(\theta\)?
\item[(ii)] if not, what is the deviation between nominal and actual
miscoverage?
\item[(iii)] what procedure can recalibrate the uncertainty set, or
stabilize the selection, so that~\eqref{eq:validity} is recovered?
\end{enumerate}

\paragraph{Strategy emphasis.}
One intended line of attack is data-driven policy selection
(\eg\ predict-then-optimize style), without classical recalibration:
randomize the data / inject structured noise---depending on inverse
optimization structure---to achieve a validity guarantee, thereby
answering~(iii).
Related touchstones include winner's-curse phenomena and algorithm
stability in post-selection inference.

\section{Part I: Stability of Constrained Optimizations}
\subsection{Problem Formulation}
Fix an optimization of the form~\eqref{eq:opt}, with data-dependent
objective \(\hat f\) and constraint \(\hat g\) learned from \(\mathcal{D}\).

The stability question is:
\begin{quote}
If the dataset changes by one data point, how much does the optimizer
output \(\hat\theta(\mathcal{D})\) change?
\end{quote}
More generally, under a specified neighbor relation on datasets
(leave-one-out, replace-one, or another structured perturbation),
bound the sensitivity of
\[
\mathcal{D}\mapsto\hat\theta(\mathcal{D})
\]
and of the induced policy / selected object.

Structural assumptions may be required on how \(\hat f\) and \(\hat g\)
are learned from \(\mathcal{D}\), and on the geometry of~\eqref{eq:opt}
(convexity, uniqueness, margin / constraint qualification, etc.).

The deliverable of Part~I is a certificate of the form
``under assumptions \(\mathcal{A}\), the map \(\mathcal{D}\mapsto\hat\theta(\mathcal{D})\)
(or the induced policy) is stable in sense \(\mathcal{S}\) with
quantitative bound \(B\).''

\section{Part II: Post-Hoc Inference after Stable Selection}
\subsection{Problem Formulation}
Assume Part~I provides a stability certificate for a (possibly
randomized) selection rule \(\widetilde A\) with output
\(\widetilde S=\widetilde A(D;\xi)\).

The inference question is:
\begin{quote}
Given that stability certificate, construct a confidence set
\(C_{\widetilde S}(D)\) such that
\[
\mathbb{P}\bigl(\theta_{\widetilde S}\in C_{\widetilde S}(D)\bigr)
\ge 1-\alpha,
\]
or, in the UQ specialization, recover~\eqref{eq:validity} when
\(\theta=\widetilde\theta(D;\xi)\) is the stabilized optimizer output.
\end{quote}

Part~II should make precise:
\begin{itemize}
\item what parameter \(\theta_{\widetilde S}\) is being covered;
\item what form of stability (\eg\ replacement stability, TV / Hellinger
stability of the selection law, score-perturbation stability) is
sufficient for the target coverage;
\item how the confidence set or uncertainty set is built from \(D\) and
\(\widetilde S\) without a naive data split, when possible.
\end{itemize}

\section{Program Decomposition}
\begin{center}
\begin{tabular}{@{}ll@{}}
\toprule
Part & Focus \\
\midrule
I & Stability of constrained data-driven optimizations \\
II & Post-hoc inference / UQ validity after stable selection \\
\bottomrule
\end{tabular}
\end{center}

The problem setting may be refined if a more principal formulation
emerges, but any revision should preserve the two-part split:
optimization stability first, then post-hoc inference on top.

\end{document}
